% =====================================================================
% demo.tex
% Minimal example exercising the definitions-latex package.
% Build with: latexmk -pdf demo.tex
% =====================================================================

\documentclass[11pt,a4paper]{article}

\usepackage[T1]{fontenc}

% Load every shortcut in one go.
\usepackage{definitions}

\title{\texttt{definitions-latex}\textemdash{}demo}
\author{}
\date{}

\begin{document}

\maketitle

\section{Differentials and partial derivatives}

\[
    \fracdif{f}{x}, \qquad
    \dift, \qquad
    \partialFrac{f}{x}, \qquad
    \parfrac{u}{t}.
\]

\section{Vectors}

\begin{description}
    \item[Cartesian basis] \(\ex \), \(\ey \), \(\ez \).
    \item[Cylindrical basis] \(\er \), \(\ep \), \(\ez \).
    \item[Spherical basis] \(\er \), \(\et \), \(\ep \).
\end{description}

\section{Brackets and norms}

\[
    \Abs{x}, \quad
    \nor{\bm{v}}, \quad
    \mean{X}, \quad
    \Par{a + b}, \quad
    \Cro{1,\, n}.
\]

\section{Integrals}

\[
    \intt f(x,y,z)\, \diftt
\]

\section{Operators}

The package defines several statistical and algebraic operators:
\section{Operators}
The package defines several statistical and algebraic operators:
\[
    \esp[X], \qquad
    \var(X), \qquad
    \sgn(x), \qquad
    \rank(M), \qquad
    \rmse, \qquad
    \pvalue.
\]
\section{Units}

Energies are written with \si{\eV}, lengths with \si{\angstrom}, pressures with
\si{\bar}, and heat with \si{\calorie}.

\section{Plasma physics shortcuts}

Some plasma-physics shortcuts are
\(\Zeff \), \(\prad \), \(\vloop \), \(\Iref \), \(\YDsC \), and \(\YOsO \).

\section{Code listings}

\begin{description}
    \item[Inline Python] \pythoncode{import numpy as np}.
    \item[Inline MATLAB] \matlabcode{x = linspace(0, 1, 100);}. % chktex 36
    \item[Inline C++] \cppcode{constexpr int n = 10;}.
\end{description}

\section{Roman numerals and the \texorpdfstring{\CC{}}{C++} logo}

Chapter \rom{4} is written in \CC{}.

\end{document}